\documentclass[10pt,a4paper]{amsart}
\usepackage[margin=19mm]{geometry}
\usepackage{amsmath,amssymb,mathtools}
\usepackage[scr=rsfs,cal=euler]{mathalfa}
\usepackage{xfrac}
\usepackage[unicode,hidelinks,bookmarks=false]{hyperref}
\newcommand{\ie}{i.e.\ }
\newcommand{\cf}{cf.\ }
\newcommand{\resp}{respectively\ }
\newcommand{\Subsection}{\subsection}
\providecommand{\nobreakdash}{\nobreak\hbox{-}}
\setlength{\emergencystretch}{2em}
\multlinegap0pt
\usepackage[shortlabels]{enumitem}
\usepackage{amssymb}
\usepackage{xcolor}
\usepackage{tikz}
\newtheorem{theorem}{Theorem}
\newtheorem{prop}{Proposition}
\newtheorem{lemma}{Lemma}
\usepackage{cleveref}
\crefname{theorem}{Theorem}{Theorems}
\crefname{prop}{Proposition}{Propositions}
\crefname{lemma}{Lemma}{Lemmas}
\newcommand{\Perm}{\mathop{\mathrm{Perm}}}
\newcommand{\ds}{\mathop{d_\square}}
\title{Chebyshev centers and radius of the set of permutons}
\author{Bal\'azs Maga}
\date{Source: arXiv:2602.02889v2 (CC BY 4.0)}
\begin{document}
\maketitle

\noindent\emph{Source and licence.} Chebyshev centers and radius of the set of permutons. Version arXiv:2602.02889v2 (CC BY 4.0), distributed by arXiv under the Creative Commons Attribution 4.0 International licence, http://creativecommons.org/licenses/by/4.0/, as recorded in the arXiv licence metadata for this exact version and shown on the abstract page https://arxiv.org/abs/2602.02889v2. Full licence notice: \texttt{licenses/arxiv-2602.02889-CC-BY-4.0.txt}.\par\smallskip

\noindent\textit{Editorial scope.} Counted unit: the article's theorem thm:main (source0.tex lines 93--95). The article's complete original proof of it is delivered with it (source0.tex lines 342--371, one \texttt{proof} environment of 30 lines, which the article places away from the statement). No mathematics is written, repaired or re-derived: the statement and the proof are the article's own lines, in the article's own notation, and no retained line is substituted or re-flowed.  Retained uncounted as the source-local context the proof needs: lines 137--145; lines 190--201.  Nothing was omitted from any retained span, so the package carries no omission record.  No retained line contains a citation, so the wrapper carries no bibliography and no bibliography entry.  No retained line contains a figure environment, an image-inclusion command or drawing code, so no image is delivered and the package declares no asset dependency.  Editorial formatting only, in addition: an AMS \texttt{amsart} preamble that restates the article's own declarations and macros used by the retained lines, namely the environments \texttt{theorem}, \texttt{prop}, \texttt{lemma} on separate counters, together with the standard packages those lines need.  The retained environments print in this document as Theorem 1, Proposition 1, Lemma 1 in order of appearance, whereas the article prints the same items with the numbering of its own full document.  The complete native source of the article as distributed in the arXiv e-print for this exact version is bundled unchanged as \texttt{sources/193870/source0.tex}.
\begin{prop} \label{prop:torus_unit_square_identification}
    For any permutons $\mu, \nu$
    \begin{enumerate}[label=(\roman*)]
        \item{For any $k, l$ integers, with the addition understood mod 2
        $$\mu(R_{i, j})-\nu(R_{i, j})=(-1)^{k+l}\left(\mu(R_{i+k, j+l})-\nu(R_{i+k, j+l})\right).$$}\label{item1:prop:torus_unit_square_identification}
        \item The value of $|\mu(R_{i, j})-\nu(R_{i, j})|$ is independent of $i, j$.
        \item $$\ds(\mu, \nu)=\max_{R\in\mathcal{R}_{\mathbb{T}^2}}|\mu(R)-\nu(R)|.$$
    \end{enumerate}
\end{prop}

\begin{lemma}\label{lemma:1/2_periodicity_equivalence}
    The following are equivalent for a permuton $\mu$ (below $w, h$ are always the sidelengths of the corresponding rectangle):
    \begin{enumerate}
        \item[(i)] $\mu$ is 1/2-periodic.
        \item[(ii)] for every toric rectangle $R$, if $h=1/2$ or $w=1/2$, then $\mu(R)=wh$.
        %\item for any standard rectangle $R$, if $h=1/2$ or $w=1/2$, then $\mu(R)=wh$.
        \item[(iii)] for every toric rectangle $R$ with $h+w\geq 1$ so that $w\leq 1/2$ or $h\leq 1/2$,
        $$\mu(R)\leq \frac{h+w}{2}-\frac{1}{4}.$$
        \item[(iv)] for every toric rectangle $R$ with $h+w=1$,
        $$\mu(R)\leq 1/4.$$
    \end{enumerate}
\end{lemma}

\begin{theorem} \label{thm:main}
    The Chebyshev radius of $\Perm$ in $\ds$ is 1/4, and a permuton is a Chebyshev center if and only if it is 1/2-periodic.
\end{theorem}

\begin{proof}[Proof of Theorem \ref{thm:main}]
    Take any permutons $\mu, \nu$ and any toric rectangle $R\in\mathcal{R}_{\mathbb{T}^2}$. Consider the corresponding toric quartet. If $\mu(R)-\nu(R)>1/2$, then by \cref{prop:torus_unit_square_identification} (i), we have $\mu(R_{1, 1})-\nu(R_{1, 1})>1/2$, which implies $\mu([0, 1]^2)>1$, an obvious contradiction. Thus the diameter is at most 1/2. Moreover, we know that $\ds(Id, 1-Id)=1/2$, as demonstrated by $R=[0, 1/2]^2$, implying that the diameter is actually 1/2.

    Concerning the Chebyshev center, as the diameter is 1/2, we know that the Chebyshev radius cannot be smaller than 1/4. 

    Assume now that $\mu$ is 1/2-periodic, and hence by Lemma \ref{lemma:1/2_periodicity_equivalence} $\mu(R)=hw$ for every rectangle with sidelengths $h, w$ if $h=1/2$ or $w=1/2$. Now it is sufficient to show that it is a Chebyshev center. Proceeding towards a contradiction, assume the existence of a permuton $\nu$ and a toric rectangle $R$ for which 
    \begin{equation}\label{eq:towards_contradiction}
    \mu(R)-\nu(R)>1/4.
    \end{equation}
    Switching to $R_{1, 1}$ if necessary, by Proposition \ref{prop:torus_unit_square_identification} (i), we can assume $h\leq 1/2$. Observe furthermore that if $w\leq 1/2$ as well, then by 1/2-periodicity $\mu(R)\leq 1/4$, ruling out this possibility. Hence $w>1/2$ can be assumed.

    Translating the measures and $R$ as well, we might assume that $R=[0, w]\times [0, h]$. Then we have
    $$R\subseteq \left([0, 1/2]\times [0, h]\right)\cup\left([1/2, w]\times [0, 1/2]\right),$$
    thus by 1/2-periodicity
    \begin{equation}\label{eq:mu_of_R}
    \mu(R)\leq \frac{h}{2}+\frac{1}{2}\left(w-\frac{1}{2}\right)=\frac{h+w}{2}-\frac{1}{4}.
    \end{equation}
    Considering \eqref{eq:towards_contradiction}, this implies $h+w> 1$.
    
    Consider now $R_{1, 0}=[w, 1]\times [0, h]$. By uniform marginals, we have
    $$\nu(R_{1, 0})\leq 1-w, \qquad \nu(R_{1, 0}\cup R)=h,$$
    thus
    $$\nu(R)\geq h-(1-w)=h+w-1.$$
    Pairing this with \eqref{eq:mu_of_R}, we obtain
    $$\mu(R)-\nu(R)\leq \frac{3}{4}-\frac{h+w}{2}<\frac{1}{4},$$
    as we already noted $h+w>1$. This ultimately contradicts \eqref{eq:towards_contradiction}. Thus indeed every $1/2$-periodic permuton $\mu$ is a Chebyshev center of $\Perm$ with respect to $\ds$.

    Finally, we prove that if $\mu$ is not 1/2-periodic then it cannot be a Chebyshev center. This follows immediately from (iv) of \cref{lemma:1/2_periodicity_equivalence}: the lack of 1/2-periodicity implies the existence of $R$ with sidelengths $w, 1-w$ and $\mu(R)>1/4$. However, for such $R$, the uniform measure $\nu$ on $R_{0, 1}\cup R_{1, 0}$ is a permuton for which $\nu(R)=0$. Thus $\ds(\mu, \nu)>1/4$.

\end{proof}

\end{document}
